\documentclass{article}

\usepackage[italian]{babel}
\usepackage[margin=20mm, footskip=20pt]{geometry}
\usepackage{graphicx}
\usepackage{subfiles}
\usepackage{hyperref}
\usepackage{nameref}
\usepackage{titlesec}
\usepackage{longtable}
\usepackage[table]{xcolor}
\usepackage{titling}
\usepackage{lastpage}
\usepackage{ifthen}
\usepackage{calc}
\usepackage{soulutf8}
\usepackage{contour}
\usepackage{float}
\usepackage{fancyhdr}
\usepackage{multirow}
\usepackage{pgfgantt}
\usepackage{lscape}
\usepackage{background}
\usepackage{lmodern}
\usepackage{textcomp}
\usepackage{lastpage}
\usepackage[utf8]{inputenc}
\usepackage{makecell}

% definizione dei percorsi in cui cercare immagini
\graphicspath{ {./}
    {./src/img/}
}

% setup della sottolineatura
\setuldepth{Flat}
\contourlength{0.8pt}

\newcommand{\uline}[1]{%
  \ul{{\phantom{#1}}}%
  \llap{\contour{white}{#1}}%
}


% setup dei link
\hypersetup{
  % set true if you want colored links (instead of boxes)
  colorlinks=true,
  % set to all if you want both sections and subsections linked
  linktoc=all,
  % set color for file links
  filecolor=blue,
  % set color for internal links
  linkcolor=black,
  % set url color
  urlcolor=blue,
  % set characters encoding in the bookmarks tab
  pdfencoding=unicode,
}

% setup forma \paragraph e \subparagraph
\titleformat{\paragraph}[hang]{\normalfont\normalsize\bfseries}{\theparagraph}{1em}{}
\titleformat{\subparagraph}[hang]{\normalfont\normalsize\bfseries}{\thesubparagraph}{1em}{}

% setup profondità indice di default
\setcounter{secnumdepth}{5}
\setcounter{tocdepth}{5}

\makeatletter %% non togliere, i comandi che definiscono i placeholder vanno qui
% esempio di utilizzo: \appendToGraphicspath{./img/} (un comando diverso per ogni path da includere)
% N.B.: ci DEVE essere un forward slash alla fine del path, a indicare che è una cartella.
\newcommand\appendToGraphicspath[1]{%
  \g@addto@macro\Ginput@path{{#1}}%
}

\newcommand{\setTitle}[1]{%
  \newcommand{\@placeholderTitle}{#1}%
}
\newcommand{\placeholderTitle}{\@placeholderTitle}

\newcommand{\setUso}[1]{%
  \newcommand{\@uso}{#1}%
}
\newcommand{\uso}{\@uso}

\newcommand{\setVersione}[1]{%
  \newcommand{\@versione}{#1}%
}
\newcommand{\versione}{\@versione}

\newcommand{\disabilitaVersione}{%
  \renewcommand{\setVersione}[1]{}%
  \renewcommand{\versione}{DISABILITATA}
}

\newcommand{\setResponsabile}[1]{%
  \newcommand{\@responsabile}{#1}%
}
\newcommand{\responsabile}{\@responsabile}

\newcommand{\setRedattori}[1]{%
  \newcommand{\@redattori}{#1}%
}
\newcommand{\redattori}{\@redattori}

\newcommand{\setVerificatori}[1]{%
  \newcommand{\@verificatori}{#1}%
}
\newcommand{\verificatori}{\@verificatori}

\newcommand{\setDestinatari}[1]{%
  \newcommand{\@destinatari}{#1}%
}
\newcommand{\destinatari}{\@destinatari}

\newcommand{\setDescrizione}[1]{%
  \newcommand{\@descrizione}{#1}%
}
\newcommand{\descrizione}{\@descrizione}

\newcommand{\setModifiche}[1]{%
  \newcommand{\@modifiche}{#1}%
}

\newcommand{\modifiche}{\@modifiche}
\makeatother %% non togliere, i comandi che definiscono i placeholder vanno qui

% hook per lo script che genera il glossario
\newcommand{\glossario}[1]{#1\textsubscript{G}}

% comandi per rendere opzionali gli elenchi di figure
\newcommand{\elencoFigure}{%
  \renewcommand{\listfigurename}{Elenco delle figure}%
  \listoffigures%
}

\newcommand{\disabilitaElencoFigure}{%
  \renewcommand{\elencoFigure}{}%
}

% comandi per rendere opzionali le tabelle
\newcommand{\elencoTabelle}{%
  \renewcommand{\listtablename}{Elenco delle tabelle}%
  \listoftables%
}

\newcommand{\disabilitaElencoTabelle}{%
  \renewcommand{\elencoTabelle}{}%
}

% comando per riferirsi ad una sezione
\newcommand{\refSec}[1]{%
  \hyperref[#1]{\S\ref*{#1}}%
}


\appendToGraphicspath{../../config/src/img/}

\setTitle{Verbale Interno 2021-04-01}

\setVersione{v0.1.0}

\setResponsabile{
	Kostadinov Samuel
}

\setRedattori{
	Paparazzo Giorgia
}

\setVerificatori{
	Rizzo Stefano
}

\setUso{Interno}

\setDestinatari{
	prof. Vardanega Tullio \\ &
	prof. Cardin Riccardo \\ &
	SpaghettiCode
}

\setDescrizione{Riassunto dell'incontro tenutosi il 1 Aprile 2021 realizzato dal gruppo SpaghettiCode in forma di meeting online.}

\setModifiche{
	v0.1.0 & Paparazzo Giorgia & Progettista & 2021-04-01 & Creazione del documento e prima stesura & Rizzo Stefano
}

\disabilitaElencoFigure
\disabilitaElencoTabelle

\begin{document}

\pagenumbering{gobble}


% variabile prima pagina (serve nell'istruzione dopo per stampare il bg solo nella prima pagina)
\newif\iffirstpage
\firstpagetrue

% setta l'immagina di bg
\backgroundsetup{
	scale=1,
	opacity=0.2,
	placement=top,
	contents={%
			\iffirstpage
				\includegraphics[width=\paperwidth]{datascience_og_colori.png}%
				\global\firstpagefalse
			\fi
		}%
}

\begin{titlepage}% per non stampare il numero della pagina

	\centering % allinea al centro la pagina
	\hspace{0.05\textwidth}% spazio tra linea e testo
	% lasciare questa riga per il corretto funzionamento di \parbox
	\parbox[b]{0.4\textwidth}{% cambiando la larghezza del testo il paragrafo si muove a destra o a sinistra 
	{\hspace{0.05\textwidth}\includegraphics[width=4cm,height=4cm]{logo_colori.png}}\\[3\baselineskip] % logo
	{\Huge\bfseries SpaghettiCode}\\ [\baselineskip] %titolo
	{\texttt{spaghetti.code.g6@gmail.com}}\\[\baselineskip]\\[4\baselineskip] % 
	{\Large\textsc\mbox{\placeholderTitle{}}}\\[4\baselineskip] % nome del documento
	{\begin{tabular}{r|l}
		\hline                                  \\
		% testo in grassetto
		\textbf{Versione}     & \versione{}     \\
		\rule{0pt}{3ex}%  EXTRA vertical height 
		\textbf{Approvazione} & \responsabile{} \\
		\rule{0pt}{3ex}%  EXTRA vertical height 
		\textbf{Redazione}    & \redattori{}    \\
		\rule{0pt}{3ex}%  EXTRA vertical height 
		\textbf{Verifica}     & \verificatori{} \\
		\rule{0pt}{3ex}%  EXTRA vertical height 
		\textbf{Uso}          & \uso{}          \\
		\rule{0pt}{3ex}%  EXTRA vertical height 
		\textbf{Destinato a}  & \destinatari{}
		\ifthenelse{\equal{\uso}{Esterno}}{
		\\ & Zucchetti S.p.A.
		}{}
	\end{tabular}}\\[4\baselineskip]

	}

	{\bfseries Descrizione}\\
	{\descrizione{}}\\[1\baselineskip]



\end{titlepage}

\newgeometry{textheight=660pt, lmargin=2cm, tmargin=2cm, rmargin=2cm}

% setup di header e footer nelle pagine senza numero
\fancypagestyle{nopage}{%
	\fancyhf{}%
	\fancyhead[R]{\includegraphics[width=1.3cm]{logo_colori.png}}%
	\fancyhead[L]{\emph{SpaghettiCode}\\\placeholderTitle{}}%
}
% setup di header e footer nelle pagine col numero
\fancypagestyle{usual}{%
	\fancyhf{}%
	\fancyhead[R]{\includegraphics[width=1.3cm]{logo_colori.png}}%
	\fancyhead[L]{\emph{SpaghettiCode}\\\placeholderTitle{}}%
	\fancyfoot[R]{\thepage\ di~\pageref{LastPage}}%
}
\setlength{\headheight}{1.8cm}

\newpage
\pagestyle{nopage}

\setcounter{table}{-1}

%REGISTRO DELLE MODIFICHE

\section*{Registro delle modifiche}
\label{sec:registro_delle_modifiche}

\rowcolors{2}{white!80!lightgray!90}{white}
\renewcommand{\arraystretch}{2} % allarga le righe con dello spazio sotto e sopra

\begin{longtable}
	[H]{|>{\centering\bfseries}m{2cm}|>{\centering}m{3.5cm}|>{\centering}m{2.5cm}|>{\centering}m{3cm}|>{\centering\arraybackslash}m{5cm}|}
	
	\hline
	\rowcolor{lightgray}
	{\textbf{Versione}} & {\textbf{Nominativo}} & {\textbf{Ruolo}} & {\textbf{Data}} & {\textbf{Descrizione}} \\
	\hline
	\endfirsthead
	
	\hline
	\rowcolor{lightgray}
	{\textbf{Versione}} & {\textbf{Nominativo}} & {\textbf{Ruolo}} & {\textbf{Data}} & {\textbf{Descrizione}} \\
	\hline
	\endhead
	
	\hline
	\rowcolor{white}
 	\multicolumn{5}{|c|}{\emph{Continua alla pagina successiva...}}\\
	\hline

	\endfoot
	\hline
	\endlastfoot
	
	\modifiche{}

\end{longtable}
% section registro_delle_modifiche (end)

\newpage
\thispagestyle{nopage}
\pagenumbering{roman}
\tableofcontents

\elencoFigure{}%
\elencoTabelle{}%

\newpage
\pagestyle{usual}
\pagenumbering{arabic}


\section{Informazioni generali}
\label{sec:info_generali}

\subsection{Informazioni incontro}
\label{sub:info_incontro}

\begin{itemize}
	\item \textbf{Luogo}: Applicazione \glossario{Discord};
	\item \textbf{Data}: 2021-04-01;
	\item \textbf{Ora}: 18:00-19.30;
	\item \textbf{Partecipanti}:
	\begin{itemize}
		\item Contro Daniel Eduardo
		\item Fichera Jacopo
		\item Kostadinov Samuel
		\item Pagotto Manuel
		\item Paparazzo Giorgia
		\item Masevski Martin
		\item Rizzo Stefano
	\end{itemize}
\end{itemize}

\section{Ordine del giorno}%
\label{sec:ordine_del_giorno}
Di seguito vengono riportati i punti discussi durante l'incontro:
\begin{enumerate}
	\item \nameref{itm:1};
	\item \nameref{itm:2};
	\item \nameref{itm:3};
	\item \nameref{itm:4};
	\item \nameref{itm:5};
    \item \nameref{itm:6};
    \item \nameref{itm:7};
    \item \nameref{itm:8};
\end{enumerate}


\subsection{Riferimenti}%
\label{sub:riferimenti}
\begin{itemize}
    \item \textsc{Piano di Progetto} v.1.4.1;
    \item \textsc{Piano di Qualifica} v.1.1.3;
    \item \textsc{Analisi dei Requisiti} v.1.3.0;
    \item \textsc{Norme di Progetto} v.2.3.0;
    \item \textsc{Glossario} v.1.2.0;
    \end{itemize}

\newpage
\section{Resoconto}
\label{sec:resoconto}

Di seguito vengono riportate le decisioni concordate e gli esiti delle diverse discussioni con il proponente, oltre a
delle note aggiuntive.

\subsection{Modifiche da apportare al Versionamento}
\label{itm:1}

Si discute una nuova modalità di versionamento da applicare ai documenti. Si decide quindi di applicare il \glossario{ versionamento semantico} e utilizzare le seguenti regole:
\begin{itemize}
    \item \textbf{X}: incrementa quando viene effettuata una \emph{modifica non retrocompatibile};
	\item \textbf{Y}: incrementa quando si effettuano \emph{aggiunte  retrocompatibili};
	\item \textbf{Z}: incrementa quando vengono effettuate \emph{modifiche retrocompatibili}.
\end{itemize}
Inoltre, ad ogni scatto di versione deve corrispondere il relativo verificatore che ha controllato la bontà della modifica.
Dal momento che il progetto sta entrando nel vivo della programmazione si decide anche di utilizzare una notazione che identifichi lo stato di avanzamento del software, si sceglie il seguente modo:
\begin{center}
    \textbf{[A].[B]}
\end{center}
dove:
\begin{itemize}
\item \textbf{B} scatterà di versione ogni volta che verrà implementato un incremento pianificato nel \textsc{Piano di Progetto};
\item \textbf{A} scatterà ogni volta che il software verrà testato e ritenuto un prodotto stabile e rilasciabile.
\end{itemize}
Si ritiene utile anche scrivere eventualmente una mail al Prof. Vardanega per capire se questo nuovo versionamento possa essere più appropriato.
\subsubsection*{Documenti coinvolti}
Tali modifiche saranno individuate fin da subito in tutti i documenti aggiornati.

\subsection{Modello di sviluppo}
\label{itm:2}

A seguito degli errori segnalati circa l'applicazione del modello di sviluppo incrementale, si vede la necessità di rivedere tale modello e adeguare la pianificazione rendendola
più orientata agli obiettivi. Verrà aggiornata la pianificazione futura esposta nel \textsc{Piano di Progetto}.
\subsubsection*{Documenti coinvolti}
Tali modifiche saranno individuabili nel \textsc{Piano di Progetto} a partire dalla versione v.1.4.1.

\subsection{Consuntivo di periodo}
\label{itm:3}

Si vede la necessità di sistemare la pianificazione futura per renderla più realistica e conforme alle esigenze maturate e all’esperienza fatta.
\subsubsection*{Documenti coinvolti}
Tali modifiche saranno individuabili nel \textsc{Piano di Progetto} a partire dalla versione v.1.4.1.

\subsection{Cruscotto di valutazione}
\label{itm:4}

Poiché il cruscotto di valutazione risulta debole si ampliano le metriche presenti, si fa corrispondere ad ogni obiettivo una metrica, inoltre
si mettono insieme i dati posseduti per avere una visione d’insieme sull’andamento complessivo.
\subsubsection*{Documenti coinvolti}
Tali modifiche saranno individuabili nel \textsc{Piano di Qualifica} a partire dalla versione v.1.1.3.

\subsection{Struttura dei verbali}
\label{itm:5}

Poiché i verbali risultano inadeguati, si decide di modificarli e rendere più strutturato il tracciamento delle decisioni. Nello specifico, infatti, si decide di inserire in ogni verbale una tabella che riassuma le decisioni prese e le identifichi con un codice.
\subsubsection*{Documenti coinvolti}
Tali modifiche saranno individuabili nei \textsc{Verbali interni}.

\subsection{Quantità dei verbali esterni}
\label{itm:6}

Poiché viene fatto notare che i verbali esterni risultano insufficienti, si decide di fare almeno un incontro con il proponente in modo da coinvolgerlo e aggiornarlo circa lo sviluppo del prodotto.
\subsubsection*{Documenti coinvolti}
Tali modifiche saranno individuabili nei Verbali esterni.

\subsection{Norme di Progetto}
\label{itm:7}

A seguito degli errori segnalati, gli \emph{amministratori} approfondiranno lo standard ISO al fine di avere maggiore consapevolezza del concetto di processo e attività. Verrà di conseguenza
effettuata una revisione complessiva delle Norme di Progetto.
\subsubsection*{Documenti coinvolti}
Tali modifiche saranno individuabili nelle \textsc{Norme di progetto} a partire dalla versione v.2.3.0.

\subsection{Indicizzazione errata nel Glossario}
\label{itm:7}

A seguito degli errori segnalati si corregge l'indicizzazione del \emph{Glossario}.
\subsubsection*{Documenti coinvolti}
Tali modifiche saranno individuabili nel \textsc{Glossario} a partire dalla versione v.1.2.0.

\subsection{Analisi dei Requisiti}
\label{itm:8}

A seguito degli errori segnalati si correggono le precisazioni fatte circa l'\textsc{Analisi dei Requisiti}.
\subsubsection*{Documenti coinvolti}
Tali modifiche saranno individuabili nell'\textsc{Analisi dei Requisiti} a partire dalla versione v.1.3.0.

\newpage
\section{Tabella delle decisioni}%
\label{sub:decisioni}

\begin{table}[!ht]
	\centering
	\begin{tabular}{|c|p{13cm}|}
		\hline
		\rowcolor{lightgray}
		\textbf{Codice} & \textbf{Descrizione} \\
		\hline
			VI0401-01 & Il versionamento verrà fatto secondo lo standard del versionamento semantico, inoltre si aggiungerà un codice di versione software. \\
			VI0401-02 & Si approfondisce il modello di sviluppo e si aggiorna la pianificazione futura basandola su obiettivi. \\
            VI0401-03 & Si adegua la pianificazione futura in base alle necessità emerse. \\
            VI0401-04 & Si amplia il cruscotto di valutazione associando ad ogni obiettivo una metrica; si aggiungono grafici che diano una visione d'insieme di andamento. \\
            VI0401-05 & Si aggiunge una tabella riepilogativa nei verbali interni. \\
            VI0401-06 & Si intrattiene un rapporto più attivo con il proponente. \\
            VI0401-07 & Si migliorano le norme approfondendo lo standard ISO. \\
            VI0401-08 & Si correggono gli errori segnalati per l'Analisi dei requisiti.  \\
		\hline
	\end{tabular}
	\caption{Tabella delle decisioni}
\end{table}

\end{document}
